\note{The first case study investigates why BBR scored so poorly under the throughput-only objective ($t_{\mathrm{coeff}}=1$). Examining the logs, the authors found that AdvNet generated an environment that maintains a high RTT for about one second and then abruptly drops the RTT. As a result, packets sent earlier during the high-RTT phase are still in flight when packets sent later during the low-RTT phase arrive at the receiver first.\par \par Using eBPF and Wireshark, the authors traced the exact sequence of events:\par 1. Seeing out-of-order packets, the receiver sends duplicate ACKs (dup ACKs), which quickly reach the sender over the now-low RTT.\par 2. Upon receiving three dup ACKs, the sender initiates a fast retransmit for the delayed packet—even though it was never lost, just delayed.\par 3. Crucially, the original delayed packet finally arrives at the receiver *just before* its fast-retransmitted copy. The receiver processes the original packet and advances its ACK number.\par 4. Immediately afterward, the fast-retransmitted copy arrives. This specific arrival order triggers a corner-case bug in the Linux kernel's ACK-update logic: even as subsequent packets arrive normally, the receiver fails to advance the acknowledgment number and keeps sending dup ACKs, eventually forcing the sender into a retransmission timeout (RTO).
\par\smallskip Source: AdvNet paper §4.4.1}
